\documentclass[11pt]{article}
\usepackage[letterpaper,margin=1in]{geometry}
\usepackage{amsmath,amssymb}
\usepackage{fancyhdr}
\usepackage{booktabs}
\usepackage{array}
\usepackage{enumitem}
\usepackage{titlesec}
\usepackage{hyperref}
\usepackage{xcolor}

\hypersetup{colorlinks=true,linkcolor=black,urlcolor=blue,citecolor=black}

\titleformat{\section}{\normalfont\Large\bfseries}{\thesection}{1em}{}
\titleformat{\subsection}{\normalfont\large\bfseries}{\thesubsection}{1em}{}

\pagestyle{fancy}
\fancyhf{}
\fancyhead[L]{Edge-Compute Pacing \& Feedback Node --- Hardware Pitch}
\fancyfoot[C]{ECE Senior Project \thepage}
\renewcommand{\headrulewidth}{0.4pt}
\renewcommand{\footrulewidth}{0pt}

\setlength{\parindent}{0pt}
\setlength{\parskip}{0.6em}

\newcolumntype{L}[1]{>{\raggedright\arraybackslash}p{#1}}

\begin{document}

\begin{center}
{\LARGE \bfseries Edge-Compute Pacing \& Feedback Node}\\[4pt]
{\large Hardware, Mechanical \& Firmware Design Pitch}\\[10pt]
\textbf{Students:} Sean Balbale \& Aidan Cullinan\\
\textbf{Program:} B.S. Electrical \& Computer Engineering, Trinity College (Class of 2027)\\
\textbf{Advisor:} Prof.\ Byers and Prof.\ Gao (TBD)\\
\textbf{Companion CS Effort:} Sean Balbale, Kayla Walsh \& Zoe Davidow, CPSC 498, advised by Prof.\ Kenneth A.\ Kousen\\
\textbf{Date:} September 2026
\end{center}

\section*{Abstract}
This is the ECE half of the Edge-Compute Pacing \& Feedback Node --- a battery-powered device that mounts to a bike or rowing shell, pairs with a wireless heart-rate strap, and gives fatigue feedback through an LED array and buzzer, offline. CPSC 498 (Prof.\ Kousen) covers the fatigue-model algorithm. This half covers microcontroller and sensor-protocol selection, power budgeting, enclosure CAD, and environmental sealing.

\section{Problem Statement}
Package a battery-powered MCU node --- calibration UI, wireless sensor link, LED/buzzer output --- into an enclosure that survives splash exposure on a shell and vibration on a bike, holding sub-100\,ms sensor-to-output latency and multi-hour runtime.

Three sub-problems:
\begin{enumerate}[topsep=2pt,itemsep=2pt]
    \item Component selection and power budgeting.
    \item RF/packaging co-design for a BLE link on a carbon or aluminum frame.
    \item Sealing an enclosure with more penetrations (buttons, OLED window, USB-C, LED window) than a typical sensor puck, under active spray.
\end{enumerate}

\section{Hardware Architecture}

\subsection{Microcontroller}
Nordic nRF52840 (Adafruit Feather nRF52840 Express). BLE~5 and ANT+ on one die, Cortex-M4F @ 64\,MHz, 256\,kB RAM, 1\,MB flash. On-board USB and LiPo charge circuitry removes a discrete charger IC.

\subsection{Sensor Protocol: BLE vs.\ ANT+}
\textbf{Decision: BLE-only for Phase 1}, ANT+ (S340 SoftDevice) as a documented fallback. Polar H10, Garmin HRM-Dual/Pro, and Wahoo TICKR all broadcast both simultaneously, so BLE covers the team's straps without an SDK migration off Arduino. ANT+'s 15 channels stay on the table for a crowded regatta start.

\subsection{Power System}
3.7\,V LiPo pouch cell, JST-PH, 2000\,mAh --- plugs directly into the Feather's charge circuit.

\begin{center}
\begin{tabular}{L{4.6cm} c c}
\toprule
\textbf{Subsystem} & \textbf{Avg.\ current} & \textbf{Duty} \\
\midrule
nRF52840 core & 5\,mA & continuous \\
BLE (1\,Hz HR notifications) & 1.5\,mA & continuous \\
NeoPixel strip (8\,px) & 15\,mA & continuous \\
OLED & 15\,mA & intermittent \\
Buzzer & 30\,mA & intermittent \\
\midrule
\textbf{Estimated avg.\ draw} & \multicolumn{2}{l}{$\sim$23\,mA} \\
\textbf{Estimated runtime} & \multicolumn{2}{l}{$\sim$85\,h vs.\ a 6\,h target} \\
\bottomrule
\end{tabular}
\end{center}

Datasheet estimates, not bench-measured. Given the margin, a 500--1200\,mAh cell is worth evaluating once real draw is measured.

\subsection{User Interface Hardware}
8-px WS2812B NeoPixel stick + piezo buzzer for output. Three 6\,mm tactile buttons (Up/Down/Select) + 0.91" 128$\times$32 I2C OLED (STEMMA QT) for calibration and sensor pairing. Config menu: \emph{Pair Sensor} $\to$ LTHR $\to$ HR$_{\text{max}}$ $\to$ HR$_{\text{rest}}$ $\to$ save/exit. Buttons are disabled mid-session.

\subsection{USB-C}
USB Type-C breakout, piggyback-soldered to the Feather's Micro-USB pads (four 30\,AWG wires to VBUS/D+/D$-$/GND). Native USB already does charging and DFU; only the connector needs adapting. Continuity bench-verified before enclosure assembly, since the joints aren't inspectable once sealed.

\section{Mechanical Design \& Waterproofing}
Handlebar mount (cycling) and hull-mount (rowing shell) share one enclosure. Rowing adds splash/spray exposure cycling doesn't have.

\begin{itemize}[topsep=2pt,itemsep=2pt]
    \item \textbf{CAD:} parametric (if possible), so PCB and battery changes don't force a rebuild. Print in ASA or PETG for UV/thermal stability outdoors.
    \item \textbf{Seam:} compressible gasket (silicone cord or printed TPU) in a channel, compressed by case screws.
    \item \textbf{Buttons:} inert mid-session, so a silicone membrane cap per shaft is sufficient --- no dynamic seal needed.
    \item \textbf{OLED/LED windows:} bonded or gasketed acrylic/polycarbonate.
    \item \textbf{USB-C:} panel-mount connector with integrated gasket, or a hinged silicone cover.
    \item \textbf{Target:} IPX4--IPX6-equivalent splash/spray resistance. No full submersion requirement --- device sits above the waterline in both mounts.
    \item \textbf{Moisture:} internal desiccant packet; consider a Gore-style pressure-equalizing vent.
\end{itemize}

Before the Phase 4 field test: directed spray at seams and each penetration, dry-out, functionality check, repeated across a few wet/dry cycles.

\section{Firmware}
FreeRTOS, four tasks: sensor (BLE HR notifications), compute (fatigue-model tick), output (NeoPixel/buzzer), UI (buttons + OLED). UI task is lowest priority and never blocks the sensor-to-output path, including during a pairing scan.

\textbf{Pairing} is user-driven: entering \emph{Pair Sensor} scans for HR Service (UUID 0x180D), lists every strap heard on the OLED, Up/Down scroll, Select confirms and stores the MAC in flash. Later boots connect directly to the stored address, falling back to a fresh scan on timeout or explicit re-pair. Replaces an earlier auto-connect-to-first-match design, which risked pairing to a teammate's strap at a crowded start line.

Core sleeps between 1\,Hz compute ticks rather than busy-waiting; validating that under real BLE traffic is a Phase 1 task.

\section{ECE Questions}
\begin{enumerate}[topsep=2pt,itemsep=3pt]
    \item Does measured current draw match the $\sim$23\,mA estimate, with margin to downsize the battery?
    \item Does BLE stay stable mounted on a carbon or aluminum frame, and does antenna placement need adjustment?
    \item With the wireless link in the loop, where does the 100\,ms latency budget actually go?
    \item Does the sealing strategy hold across wet/dry cycles and vibration, or does one penetration (buttons, USB-C, window) turn out to be the weak point?
    \item How does the antenna perform with the enclosure printed in ASA or PETG, and does the material choice need to change/do we need to make a new antenna?
\end{enumerate}

\section{Bill of Materials}
\begin{center}
\begin{tabular}{L{6.2cm} c L{5.5cm}}
\toprule
\textbf{Item} & \textbf{Price} & \textbf{Notes} \\
\midrule
Adafruit Feather nRF52840 Express & \$24.95 & Core MCU \\
LiPo 3.7\,V 2000\,mAh, JST-PH & \$12.50 & \\
NeoPixel Stick, 8$\times$WS2812/SK6812 & \$5.95 & \\
Buzzer 5V & \$1.95 & \\
Tactile Button Switch (6\,mm), 20 pk & \$4.95 & 3 used \\
0.91" 128$\times$32 I2C OLED, STEMMA QT & $\sim$\$15 & \\
STEMMA QT to Header Cable & \$0.95 & \\
USB-C Breakout Board & \$2.95 & \\
Protoboard, wire, JST-PH connectors (local) & $\sim$\$18 & \\
\midrule
\textbf{Subtotal} & $\sim$\textbf{\$69} & \\
\bottomrule
\end{tabular}
\end{center}

\section{Timeline}
\begin{center}
\begin{tabular}{c L{9.5cm} c}
\toprule
\textbf{Phase} & \textbf{Milestone} & \textbf{Target} \\
\midrule
1 & Hardware bring-up; BLE pairing verified; power budget bench-validated & Sep 2026 \\
2 & Enclosure CAD v1 dry-fit; sealing approach prototyped & Oct 2026 \\
3 & FreeRTOS integration; latency/scheduling benchmarks & Nov 2026 \\
4 & Sealed enclosure passes wet/dry bench test; on-bike/on-water field test & Dec 2026 \\
5 & Write-up, poster, demo & Apr 2027 \\
\bottomrule
\end{tabular}
\end{center}

\section{Open Risks}
\begin{itemize}[topsep=2pt,itemsep=2pt]
    \item ANT+ needed only if a non-dual-mode strap turns up.
    \item Sealing approach needs bench validation before the Phase 4 field test.
    \item Battery downsize waits on real current-draw data.
    \item Handlebar mount: reuse a standard bike-computer pattern, or design a custom clamp.
\end{itemize}

\section{What We're Asking of an Advisor}
\begin{itemize}[topsep=2pt,itemsep=2pt]
    \item Review of the MCU, power, and protocol decisions as they move from datasheet estimates to bench numbers.
    \item Guidance on sealing multiple penetrations on a 3D-printed enclosure.
    \item Feedback on antenna placement and enclosure material once mounted on a carbon frame or shell.
    \item Help setting up a credible splash/wet-dry bench test.
\end{itemize}

\vspace{6pt}
\textbf{Contact:} Sean Balbale --- \href{mailto:sean.balbale@trincoll.edu}{sean.balbale@trincoll.edu}\\
Aidan Cullinan --- \href{mailto:aidan.cullinan@trincoll.edu}{aidan.cullinan@trincoll.edu}

\end{document}